\documentclass[12pt]{article}

% ---- Packages ----
\usepackage[T1]{fontenc}


\usepackage{geometry}
\geometry{a4paper, margin=1.2in}
\usepackage{amsmath}
\usepackage{amssymb}
\usepackage{amsthm}
\usepackage{hyperref}
\usepackage{csquotes}
\usepackage{enumitem}
\usepackage[nameinlink,capitalise]{cleveref}
\usepackage{natbib}
\bibliographystyle{plainnat}
\setcitestyle{authoryear,open={(},close={)}}

% ---- Theorem environments ----
\newtheorem{theorem}{Theorem}[section]
\newtheorem{corollary}{Corollary}[theorem]
\newtheorem{proposition}{Proposition}[section]
\newtheorem{definition}{Definition}[section]
\newtheorem{remark}{Remark}[section]

% ---- Macros ----
\newcommand{\bind}{\phi}
\newcommand{\recomp}{\rho}
\newcommand{\Jsel}{J_{\mathrm{sel}}}
\newcommand{\Jclos}{J_{\mathrm{clos}}}
\newcommand{\Jgrowth}{J_{\mathrm{growth}}}
\newcommand{\dsem}{d_{\mathrm{sem}}}

% ---- Formatting ----
\setlength{\parskip}{0.5em}
\setlength{\parindent}{1.5em}
\sloppy

% ---- Title ----
\title{Constraint Regimes and the Structure of Surviving Objects\\
\large On the Invariance of Selection Pressure Across Censorship, Feed Dynamics, and Cognitive Strategy}
\author{Flyxion}
\date{}

\begin{document}

\maketitle

\begin{abstract}
Three phenomena that appear historically and culturally unrelated—the formal
properties of Soviet-era science fiction, the structural dynamics of feed-optimized
media platforms, and the cognitive trajectory of agents who voluntarily adopt
high-discipline intellectual regimes—are shown to be instances of a single
underlying mechanism. In each case, a selection environment imposes a
fragment-level survivability condition, and the global form of what survives is
determined by that condition rather than by the intentions of the agents
involved. We formalize this through the coupled variables of binding force
$\bind(c)$—the degree to which a piece of content retains semantic identity
under arbitrary decomposition—and recomposability $\recomp(c)$—the degree
to which its affective or functional charge survives context-switching. Systems
that select for high $\recomp$ converge to structurally shallow, fragment-safe
outputs. Systems that select for high $\bind$ converge to internally coherent
but low-circulation objects. Neither outcome depends on stated intent. The
invariant is the constraint structure. As a corollary, voluntary adoption of a
high-binding regime functions as an exit from recomposability-dominated
selection environments, and the colloquial description of this strategy as
``galaxy brain'' discipline masks a formal structural shift. We conclude that
the full version of this claim cannot itself propagate without transformation,
which is not a paradox but a fixed point of the framework.
\end{abstract}

\tableofcontents
\clearpage

% ============================================================
\section{Introduction: Three Phenomena, One Mechanism}
\label{sec:intro}
% ============================================================

Consider three situations that appear to have nothing in common.

First: a pair of writers in the Soviet Union produce science fiction novels
spanning decades. Their work is formally published, widely read, and broadly
regarded as subversive. It escapes censorship not by avoiding critique but by
being structured such that every extractable fragment can be defended as
allegory, genre fiction, or ambiguity. The corpus feels ``perfectly calibrated,''
as though a sophisticated editorial intelligence had tuned each sentence to
the threshold of permissibility.

Second: feed-optimized media platforms select for content that generates
response without requiring full processing—content that can be clipped,
quoted, embedded, and reacted to without loss of affective charge. The
discourse that survives this selection is increasingly characterized by high
affective intensity and low semantic stability: it circulates widely but means
something different in every context it enters. This is not because the
platform intended to destroy meaning. It is because the platform optimized
for a property that is approximately inverse to meaning-preservation.

Third: an individual notices that their intellectual environment has progressively
supplied fewer genuinely difficult inputs—not because the world has become
less interesting, but because their social network can no longer generate
material above their internal processing threshold. To continue developing,
they shift from consuming externally curated complexity to constructing it
internally: writing formal papers, building systems, enforcing compression,
and selecting reading for maximal transfer across domains. The output of this
phase is structurally different from the input phase: less recomposable, more
coherent, lower in circulation, higher in internal consistency.

These three situations look unrelated. One is about political censorship, one
about media economics, one about personal intellectual development. But
once they are expressed in terms of constraint structure rather than domain
content, they become instances of the same mechanism:

\begin{quote}
\emph{Any environment that imposes selection on fragments will shape the
internal structure of what survives. The trade-off between binding force and
recomposability determines whether surviving structures remain deep or
become shallow.}
\end{quote}

This paper develops that claim formally, traces it through each of the three
cases, and examines what follows from their unification.

% ============================================================
\section{The Core Variables}
\label{sec:variables}
% ============================================================

\begin{definition}[Binding Force]
\label{def:binding}
The binding force of a content object $c$, denoted $\bind(c) \in [0,1]$,
measures the degree to which the semantic content of $c$ is preserved under
arbitrary context transformation. Formally, let $T_\sigma$ denote a
context-transformation operator parameterized by context $\sigma$. Then:
\[
\bind(c) = 1 - \mathbb{E}_\sigma\bigl[\dsem(c,\, T_\sigma(c))\bigr]
\]
where $\dsem$ is a semantic distance measure and the expectation is taken
over the distribution of contexts into which $c$ may be reinserted.
High $\bind(c)$ indicates that meaning is preserved under context-switching.
Low $\bind(c)$ indicates that affective charge survives while semantic content
drifts.
\end{definition}

\begin{definition}[Recomposability]
\label{def:recomp}
The recomposability of $c$, denoted $\recomp(c) \in [0,1]$, measures the
degree to which the affective or functional charge of $c$ is preserved under
context transformation:
\[
\recomp(c) = \mathbb{E}_\sigma\bigl[A(T_\sigma(c))\bigr] \,/\, A(c)
\]
where $A(c)$ is the affective charge of $c$. Content is maximally
recomposable when it generates equivalent response regardless of the
context into which it is inserted.
\end{definition}

\begin{proposition}[Approximate Inversion]
\label{prop:inversion}
Under typical platform dynamics, $\recomp(c) \approx 1 - \bind(c)$. More
precisely:
\[
\recomp(c) \leq 1 - \bind(c) + \varepsilon(c)
\]
where $\varepsilon(c) \geq 0$ measures the degree to which the affective
charge of $c$ has become decoupled from its semantic content through prior
recomposition or deliberate design. Content with $\varepsilon(c) = 0$ cannot
be recomposed without losing both affect and meaning simultaneously.
\end{proposition}

The quantity $\varepsilon(c)$ indexes what may be called \emph{pseudo-binding}:
the appearance of semantic constraint without actual resistance to
recomposition. Content that performs depth while remaining affectively
portable has high $\varepsilon(c)$. It is the dominant rhetorical mode of
feed-adapted discourse.

The inversion in \cref{prop:inversion} is the load-bearing element of the
entire framework. If increasing $\recomp$ necessarily decreases $\bind$,
then any selection environment that rewards $\recomp$ will systematically
erode $\bind$, regardless of what the agents within it intend. No one has to
decide to simplify things. The selection mechanism does it automatically.

% ============================================================
\section{Selection Regimes and Closure Regimes}
\label{sec:regimes}
% ============================================================

\begin{definition}[Selection Regime]
\label{def:selreg}
A discourse or production environment operates in the \emph{selection regime}
when its dominant selection functional rewards recomposability:
\[
\Jsel(c) = \alpha \cdot E(c) + \beta \cdot \recomp(c) - \gamma \cdot K(c)
\]
where $E(c)$ is engagement, $\recomp(c)$ is recomposability, $K(c)$ is
cognitive cost to the consumer, and $\alpha, \beta, \gamma > 0$. Content
survives by being detachable from context. The stable equilibrium is
low-binding, high-recomposability discourse.
\end{definition}

\begin{definition}[Closure Regime]
\label{def:closreg}
An environment operates in the \emph{closure regime} when its dominant
selection functional rewards semantic invariance under transformation:
\[
\Jclos(c) = \lambda \cdot \bind(c) - \mu \cdot \delta(c)
\]
where $\bind(c)$ is binding force and $\delta(c)$ is detected inconsistency
under adversarial recomposition. Content survives by remaining identical to
itself across contexts. The stable equilibrium is high-binding, low-entropy
discourse.
\end{definition}

\begin{theorem}[Regime Incompatibility]
\label{thm:incompatibility}
For any content $c$:
\[
\nabla_c \Jsel \cdot \nabla_c \Jclos \leq 0
\]
with equality only on the set of measure zero where $\bind(c) = \recomp(c)
= \tfrac{1}{2}$. Gradient ascent under $\Jsel$ is gradient descent under
$\Jclos$ and vice versa. There is no smooth deformation that transforms a
selection-regime system into a closure-regime system.
\end{theorem}

\begin{proof}[Proof sketch]
$\Jsel$ increases in $\recomp(c) \approx 1 - \bind(c)$, so $\nabla_c \Jsel$
points in the direction of increasing $\recomp$ and decreasing $\bind$.
$\Jclos$ increases in $\bind(c)$, so $\nabla_c \Jclos$ points in the
direction of increasing $\bind$ and decreasing $\recomp$. These gradients
are antiparallel.
\end{proof}

The two regimes are not points on a spectrum. They are characterized by
incompatible selection functionals, incompatible stable equilibria, and
incompatible populations of content. This is why attempts to operate within
a selection regime while preserving closure-regime properties—the ``Trojan
horse'' strategy—fail structurally rather than contingently.

\begin{corollary}[Trojan Horse Bound]
\label{cor:trojan}
Let $c$ be content with binding force $\bind(c)$. Let $c'$ be its
translation into a more recomposable form with $\recomp(c') > \recomp(c)$.
Then:
\[
\bind(c') \leq \bind(c)
\]
with equality only if the translation is semantically lossless, which
requires $\recomp(c') = \recomp(c)$. Since the translation is performed
precisely to increase $\recomp(c')$, we have $\bind(c') < \bind(c)$
strictly. The propagating form binds less than the original.
\end{corollary}

The corollary removes a comforting belief: that ideas can be made widely
transmissible without alteration. Once the trade-off is formalized,
accessibility is not neutral. It is a transformation operator that necessarily
changes the object. The ``smuggled'' idea is not the same idea; it is a
projection that survives under higher recomposability constraints.

% ============================================================
\section{Three Instances of the Same Structure}
\label{sec:instances}
% ============================================================

\subsection{Soviet-Era Fiction: Hostile Extraction as Selection Pressure}
\label{ssec:soviet}

The Strugatsky brothers published science fiction in the Soviet Union across
several decades. Their work was broadly regarded as subversive, yet it
escaped suppression. The standard explanation attributes this to authorial
craft: deliberate calibration of transgressive content to the threshold of
permissibility.

This explanation is correct in its description but wrong in its causal
structure. It implies a centralized optimizing agent—an author or a regime
monitoring discourse—producing the calibration through intent. What
actually produced it was a selection mechanism operating at the fragment
level.

The operative constraint was not: \emph{``this book as a whole must be
defensible.''} It was: \emph{``any extractable fragment of this book,
removed from context, must be defensible as fiction, allegory, or
ambiguity.''} This is a fragment-level survivability condition—exactly the
condition that maximizes $\bind(c)$ under hostile interpretation rather than
maximizes $\recomp(c)$ under feed amplification.

The result is a corpus where each piece is locally deniable under arbitrary
extraction. Every paragraph about a medieval tyrant can be read as genre
worldbuilding. Every passage about a population subjected to daily seizures
of manufactured euphoria can be read as speculative fiction. The work feels
``perfectly calibrated'' not because it was explicitly engineered that way,
but because everything that was not fragment-safe failed to survive.

\begin{remark}
The Soviet censorship regime and the feed selection regime impose
structurally identical fragment-level survivability conditions while
optimizing for opposite properties. Censorship selects for maximum
$\bind$ under hostile decomposition—every fragment must resist lethal
extraction. Feeds select for maximum $\recomp$ under algorithmic
amplification—every fragment must circulate without loss of charge.
Both produce corpora where no fragment is ``unsafe'' in the relevant
sense. Different penalty functions, same formal constraint structure,
opposite equilibria.
\end{remark}

The conspiracy version of this observation—that the Strugatskys were
managed by a state apparatus monitoring public discourse and serializing
feedback—is narratively satisfying but explanatorily redundant. The
selection mechanism produces the same output without coordination. This
is not to say the conspiracy version is false; it is to say it is a
limiting case of the selection mechanism, not the mechanism itself.

\subsection{Feed-Optimized Platforms: Recomposability as Fitness}
\label{ssec:feed}

Feed architectures optimize for a property of content that is not identical to
engagement but produces it: recomposability. A piece of content is
recomposable to the degree that it generates response without requiring the
respondent to have fully processed it. High recomposability content can be
clipped, quoted, embedded, and reacted to across contexts without loss of
affective charge.

The binding force of such content is correspondingly low. High-binding
content—content whose conclusion depends on the specific structure of its
premises, whose meaning changes when removed from context—is maximally
resistant to recomposition and therefore maximally penalized by the
selection functional $\Jsel$.

The practitioner class that has converged to high-$\recomp$ production has
not found a strategy for operating within a degraded environment. It has
converged to the only stable trajectory the environment admits. This
convergence does not require motivation or awareness. It requires only that
practitioners who do not converge lose visibility, and practitioners who do
converge retain it. Over time, the visible population is the converged
population.

\begin{theorem}[Convergence to Low-Binding Equilibrium]
\label{thm:convergence}
Under a platform selection functional that rewards recomposability, and
under boundary conditions that do not introduce strong countervailing
pressures, the visible discourse converges to a low-binding equilibrium
from which perturbations by high-binding content cannot produce
asymptotic stabilization.
\end{theorem}

\begin{proof}[Proof sketch]
By definition, $\Jsel$ assigns positive weight to $\recomp(c) \approx
1 - \bind(c)$. Content with high $\bind$ receives systematically lower
$\Jsel$ scores. Visibility is a monotone function of $\Jsel$. As
practitioners observe visibility distributions, adaptive pressure pushes
production toward low-binding strategies. A high-binding injection—a
single dense argument, a formal proof, a structurally demanding
essay—constitutes a local perturbation to the discourse state, but does
not alter the platform architecture or the boundary conditions. The
dynamics continue to be dominated by the gradients of $\Jsel$, which
re-amplify low-binding content and attenuate high-binding content. The
perturbation decays.
\end{proof}

\begin{corollary}[Insufficiency of Adaptive Strategies]
Adaptive strategies that reduce binding force to increase recomposability
do not alter the attractor structure of the system. They accelerate
convergence to the low-binding equilibrium.
\end{corollary}

This corollary reformulates the standard paradox of media critique: why do
sophisticated, well-intentioned practitioners keep producing content that
reinforces the systems they are analyzing? The standard explanation
invokes co-optation or false consciousness. The structural explanation is
simpler. The practitioners who remain visible are precisely those whose
output satisfied the recomposability constraint. The good versions did not
get suppressed in a conspiratorial sense. They failed to meet the
propagation threshold and became invisible. What looks like ideological
drift is survivorship bias over outputs.

\subsection{Voluntary Closure: The Self-Imposed Binding Regime}
\label{ssec:voluntary}

The third instance differs from the first two in that the selection pressure is
self-imposed rather than externally enforced. But the formal structure is
identical.

An individual operating within an externally curated intellectual environment
receives inputs filtered by other people's perception of difficulty. This
functions as a distributed selection mechanism: the social network routes
material above its own processing threshold toward the individual perceived
as capable of handling it. The incoming material is already filtered for
complexity. The individual's intellectual development occurs within a
propagation regime, except that the selection criterion is perceived
competence rather than recomposability.

This arrangement saturates. Once the environment can no longer generate
inputs above the individual's internal binding threshold—once nothing the
network sends constitutes genuine difficulty—the distributed intake system
collapses. The environment routes fewer hard problems not because the
world has fewer hard problems, but because the social network has
exhausted its capacity to identify problems hard enough to be useful.

Continuing development at this point requires switching regimes. The
individual must shift from consuming complexity to constructing it:
generating formal problems, enforcing compression across domains,
selecting reading for maximal transfer, producing objects with high
internal coherence. This shift increases $\bind$ at the expense of $\recomp$.
The outputs become less transmissible, more internally structured, and
more resistant to decomposition.

\begin{remark}
The saturation of the externally curated input channel is not a failure
state. It is a success metric: evidence that the first phase of
distributed intake operated correctly. The appropriate response is not
to seek out new sources of recomposable complexity but to switch
selection functionals—from $\Jsel$ to $\Jclos$, in effect. The second
phase is structurally harder because it removes the passive intake
channel. It is also the only phase that scales indefinitely, because
self-generated problems have no external ceiling.
\end{remark}

The individual in this third case is doing what the Strugatskys did
unintentionally and what the feed resists structurally: constructing objects
whose internal coherence is not contingent on context. The apparent
``galaxy brain'' quality—the tendency to see structural patterns everywhere,
to prefer domains where performance depends on skill rather than chance, to
read for transfer rather than coverage—is the behavioral signature of
operating under a self-imposed closure regime. It is not a personality trait.
It is a constraint structure.

% ============================================================
\section{The Unified Principle}
\label{sec:unified}
% ============================================================

The three instances share a single formal property. In each case:

\begin{enumerate}[label=\textbf{\arabic*.}]
\item A selection environment imposes a condition at the fragment level.
\item Objects that satisfy the condition survive; objects that do not,
      do not.
\item The global form of the surviving corpus is determined by the
      fragment-level condition rather than by the stated goals of the
      agents involved.
\item The apparent ``calibration'' of the output—its apparent fit to the
      selective environment—requires no coordinating agent. It is a
      consequence of the filter.
\end{enumerate}

Stated as a general principle:

\begin{quote}
\textbf{Constraint Invariance:} Any environment that imposes selection on
fragments will shape the internal structure of what survives. The
binding-recomposability trade-off determines whether surviving structures
remain deep or become shallow. The apparent intentions of the selecting
agent are irrelevant. The invariant is the constraint structure.
\end{quote}

This principle cuts through most of the usual explanatory apparatus in media
critique, political analysis, and accounts of intellectual development. It
does not require invoking coordination, conspiracy, manipulation, or
co-optation. It requires only specifying the selection functional and
reading off the equilibrium.

The unification also removes a structural asymmetry in how we typically
think about these cases. Soviet censorship is described as externally
imposed and coercive. Feed dynamics are described as emergent and
impersonal. Personal intellectual development is described as chosen and
autonomous. These descriptions are accurate about the surface mechanism.
They obscure the formal identity of the underlying structure. Specifying
the constraint reveals what the surface descriptions conceal.

% ============================================================
\section{Corollary: The Galaxy Brain as Regime Exit}
\label{sec:galaxybrain}
% ============================================================

The colloquial term ``galaxy brain'' originates in internet culture as a
satirical description of reasoning that escalates through increasingly
elaborate inferential steps to conclusions that appear absurd. The meme
encodes a double register: sincere admiration for the capacity to operate
at high levels of abstraction, and ironic awareness of the risk of
decoupling from empirical feedback.

This double register is structurally exact. It describes the phenomenology
of operating under a self-imposed closure regime from the perspective of
someone observing from within a selection regime. The outputs appear
``galaxy-brained''—elaborate, internally consistent, disconnected from the
propagation economy—precisely because they have been optimized for
$\bind$ rather than $\recomp$.

The irony in the meme performs the same function as the science-fiction
framing in the Strugatsky corpus: it provides protective ambiguity. ``I am
describing a cognitive strategy'' and ``I am mocking a cognitive failure''
are simultaneously available interpretations. The content survives under
hostile extraction—it can be read either way—because neither reading
requires retracting the other.

\begin{proposition}[Galaxy Brain as Closure Strategy]
\label{prop:galaxybrain}
Voluntary adoption of behaviors associated with ``galaxy brain''
discipline—preference for strategy over chance, reading for transfer,
compression of ideas into dense representations, avoidance of
recomposability-optimized content domains—functions as a self-imposed
closure regime. The colloquial framing masks a formal structural shift:
from optimizing for propagation to optimizing for coherence. The
cognitive gains are a byproduct of sustained exposure to high-binding
problem classes rather than a direct effect of claimed identity.
\end{proposition}

The practical implication is not to ``think like a galaxy brain'' in the
sense of adopting a self-flattering identity. It is to change the selection
functional governing input and output. Specifically: select inputs by
binding force rather than by recomposability, produce outputs that require
context to be evaluated rather than outputs that circulate without it, and
tolerate the reduction in visibility that results.

The resulting corpus—formal papers in obscure repositories, dense
monographs with limited circulation, systems that require significant prior
investment to understand—is not a failure mode. It is the signature of
correct operation under a self-imposed closure regime. It is also, by the
Trojan Horse Bound, the only kind of output whose binding force is
preserved rather than sacrificed in the act of production.

% ============================================================
\section{The Fixed Point}
\label{sec:fixedpoint}
% ============================================================

This paper cannot propagate in its current form without transforming into
something with lower binding force. Any version that achieves sufficient
circulation to be widely recognized will have undergone the transformation
$\bind(c') < \bind(c)$ that the Trojan Horse Bound predicts. The viral
summary will be recomposable. The bound version will be invisible.

This is not a paradox. It is a fixed point of the framework applied
recursively. The paper predicts its own fate: wide readership is mild
evidence against the integrity of the transmitted version. The correct
publication strategy follows directly: maximize $\bind$, deposit in a
context accessible to closure-regime readers, and do not attempt to
optimize for propagation. A version optimized for propagation would
refute itself by existing.

The framework does not conclude with a reform proposal. Reform proposals
that circulate within selection-regime environments have been selected for
recomposability and have therefore already had their binding force reduced.
The alternative to the selection regime already exists. It is structured by
a different selection functional, it produces a different kind of object, and
it is invisible from within the propagation economy because visibility within
the propagation economy requires the property that the closure regime, by
its structure, refuses to provide.

The full version of this claim cannot propagate without changing form. That
is not a counsel of despair. It is the most compressed statement of the
argument.

% ============================================================
\section{Conclusion}
\label{sec:conclusion}
% ============================================================

We began with three apparently unrelated phenomena and arrived at a single
mechanism. Soviet-era science fiction, feed-optimized media dynamics, and
voluntary high-discipline intellectual development are all expressions of the
same underlying structure: selection environments that impose
fragment-level survivability conditions shape the global form of surviving
objects through the binding-recomposability trade-off.

The conceptual contribution is the removal of motive from the explanatory
apparatus. Whether the selecting agent is a state censor, an algorithmic
recommendation system, or an individual imposing discipline on their own
cognitive inputs, the behavior of the resulting corpus is determined by the
constraint structure, not by intent. This replacement of intentional
explanation with structural explanation is not cynicism. It is precision.

The practical contribution is the identification of regime exit as an
available option. The selection regime is not the only stable environment
for intellectual production. The closure regime exists, operates under
different principles, and produces objects with different properties. The
transition requires accepting reduced recomposability—reduced
transmissibility, reduced visibility, reduced circulation—in exchange for
increased binding force. This is not a sacrifice. It is a change of
optimization target.

The framework eating itself—this paper predicting its own limited
propagation—is the cleanest confirmation that the mechanism is correctly
identified. An argument that predicts its own fate under the conditions it
describes is not self-undermining. It is self-consistent.

\end{document}
